% A Lean-oriented proof of the single-sum formula for Apéry's companion.
% Convention: binomial coefficients are zero outside their natural range.
\documentclass[11pt]{article}
\usepackage{amsmath,amssymb,amsthm,mathtools}

\newcommand{\Hthree}[1]{H_{#1}^{(3)}}
\newcommand{\Kern}{F}
\newcommand{\Ppoly}{P}
\newcommand{\dd}{\mathop{}\!d}
\newtheorem{lemma}{Lemma}

\begin{document}

\section*{A telescoping proof of Apéry's companion formula}

For $m\geq 0$, put
\[
  \Hthree m:=\sum_{j=1}^{m}\frac1{j^3},\qquad \Hthree0=0,
\]
and, for $n,k\in\mathbb Z$, put
\[
 \Kern_{n,k}:=
 \begin{cases}
 \displaystyle\binom nk^2\binom{n+k}{k}^2,&0\leq k\leq n,\\[1ex]
 0,&\text{otherwise}.
 \end{cases}
\]
Define
\[
 \widetilde b_n:=\sum_{k\in\mathbb Z}\Kern_{n,k}
       \bigl(2\Hthree n-\Hthree k\bigr)\qquad(n\geq0).
\]
The sum is finite by definition.  We prove that it is the usual Apéry
companion sequence, i.e.
\[
 \boxed{\quad b_n=\sum_{k=0}^n\binom nk^2\binom{n+k}{k}^2
       \bigl(2H_n^{(3)}-H_k^{(3)}\bigr).\quad}                 \tag{A}\label{eq:A}
\]

Let
\[
 \Ppoly(n):=34n^3+51n^2+27n+5,
 \qquad
 d_{n,k}:=2\Hthree n-\Hthree k.
\]
We use the following two rational identities.  They are the only
non-formal algebraic input in the proof.

\begin{lemma}[Apéry kernel certificate]\label{lem:kernel}
For $n\geq1$ and $k\in\mathbb Z$,
\[
 (n+1)^3\Kern_{n+1,k}-\Ppoly(n)\Kern_{n,k}+n^3\Kern_{n-1,k}
 =G_{n,k}-G_{n,k-1},                                          \tag{B}\label{eq:B}
\]
where
\[
 G_{n,k}:=4(2n+1)\bigl(k(2k+1)-(2n+1)^2\bigr)\Kern_{n,k}.
\]
\end{lemma}

\begin{lemma}[Residual certificate]\label{lem:residual}
For $n\geq1$ and $k\in\mathbb Z$, the expression
\[
 \frac{G_{n,k-1}}{k^3}+2\Kern_{n+1,k}-2\Kern_{n-1,k}          \tag{C}
\]
is understood as $0$ when $k=0$, and otherwise equals
\[
 C_{n,k+1}-C_{n,k},                                           \tag{D}
\]
where
\[
 C_{n,k}:=
 \frac{4k^4(k-1)(2n+1)\bigl(k-2n(n+1)-1\bigr)}{n^2(n+1)^2}
 \frac{(n+k-1)!^2}{k!^4(n-k+1)!^2}.
\]
Here $C_{n,k}$ is set to $0$ whenever one of the displayed factorials
is outside its natural range.
\end{lemma}

\begin{proof}[Proof of \eqref{eq:A} from Lemmas \ref{lem:kernel}--\ref{lem:residual}]
The elementary harmonic identities are
\[
 d_{n+1,k}-d_{n,k}=\frac{2}{(n+1)^3},\qquad
 d_{n,k}-d_{n-1,k}=\frac{2}{n^3},\qquad
 d_{n,k}-d_{n,k-1}=-\frac1{k^3}.                              \tag{E}\label{eq:E}
\]
For the last formula we only use $k\geq1$.

Multiply \eqref{eq:B} by $d_{n,k}$ and use the first two identities in
\eqref{eq:E}.  This gives
\begin{align*}
 &(n+1)^3\Kern_{n+1,k}d_{n+1,k}
       -\Ppoly(n)\Kern_{n,k}d_{n,k}
       +n^3\Kern_{n-1,k}d_{n-1,k}\\
 &\quad=(G_{n,k}-G_{n,k-1})d_{n,k}
       +2\Kern_{n+1,k}-2\Kern_{n-1,k}\\
 &\quad=G_{n,k}d_{n,k}-G_{n,k-1}d_{n,k-1}
       +\frac{G_{n,k-1}}{k^3}
       +2\Kern_{n+1,k}-2\Kern_{n-1,k}.
                                                                    \tag{F}\label{eq:F}
\end{align*}
At $k=0$ the $G_{n,-1}/k^3$ term is omitted, and the same formula holds.
By Lemma~\ref{lem:residual}, \eqref{eq:F} becomes
\begin{align*}
 &(n+1)^3\Kern_{n+1,k}d_{n+1,k}
       -\Ppoly(n)\Kern_{n,k}d_{n,k}
       +n^3\Kern_{n-1,k}d_{n-1,k}\\
 &\qquad=
 \bigl(G_{n,k}d_{n,k}+C_{n,k+1}\bigr)
 -\bigl(G_{n,k-1}d_{n,k-1}+C_{n,k}\bigr).                    \tag{G}\label{eq:G}
\end{align*}
Summing \eqref{eq:G} over $k\in\mathbb Z$ is legitimate because every
term has finite support.  The right-hand side telescopes to zero, hence
\[
 (n+1)^3\widetilde b_{n+1}-\Ppoly(n)\widetilde b_n
       +n^3\widetilde b_{n-1}=0\qquad(n\geq1).                \tag{H}\label{eq:H}
\]
Finally,
\[
 \widetilde b_0=0,
 \qquad
 \widetilde b_1=\Kern_{1,0}(2\Hthree1-\Hthree0)
          +\Kern_{1,1}(2\Hthree1-\Hthree1)=2+4=6.
\]
Thus $(\widetilde b_n)$ has Apéry's companion recurrence and its two
defining initial values $(0,6)$.  Uniqueness of a second-order recurrence
therefore gives $\widetilde b_n=b_n$ for every $n\geq0$.
\end{proof}

\paragraph{Formalisation split.}
The Lean development naturally separates into: (i) proving the two
certificates \eqref{eq:B} and \eqref{D} by field normalisation after
expanding the binomial/factorial ratios; (ii) the finite-support telescoping
argument \eqref{eq:G}$\Rightarrow$\eqref{eq:H}; and (iii) recurrence uniqueness
from the initial values.  It is usually easiest to formalise the certificates
on the finite index interval $0\leq k\leq n+1$, avoiding a globally defined
factorial quotient for $C_{n,k}$.

\end{document}
